% Options for packages loaded elsewhere
\PassOptionsToPackage{unicode}{hyperref}
\PassOptionsToPackage{hyphens}{url}
%
\documentclass[
  11pt,
  a4paper,
]{article}
\usepackage{amsmath,amssymb}
\usepackage{iftex}
\ifPDFTeX
  \usepackage[T1]{fontenc}
  \usepackage[utf8]{inputenc}
  \usepackage{textcomp} % provide euro and other symbols
\else % if luatex or xetex
  \usepackage{unicode-math} % this also loads fontspec
  \defaultfontfeatures{Scale=MatchLowercase}
  \defaultfontfeatures[\rmfamily]{Ligatures=TeX,Scale=1}
\fi
\usepackage{lmodern}
\ifPDFTeX\else
  % xetex/luatex font selection
\fi
% Use upquote if available, for straight quotes in verbatim environments
\IfFileExists{upquote.sty}{\usepackage{upquote}}{}
\IfFileExists{microtype.sty}{% use microtype if available
  \usepackage[]{microtype}
  \UseMicrotypeSet[protrusion]{basicmath} % disable protrusion for tt fonts
}{}
\makeatletter
\@ifundefined{KOMAClassName}{% if non-KOMA class
  \IfFileExists{parskip.sty}{%
    \usepackage{parskip}
  }{% else
    \setlength{\parindent}{0pt}
    \setlength{\parskip}{6pt plus 2pt minus 1pt}}
}{% if KOMA class
  \KOMAoptions{parskip=half}}
\makeatother
\usepackage{xcolor}
\usepackage[margin=3cm,marginparwidth=2cm]{geometry}
\usepackage{listings}
\lstset{defaultdialect=[5.3]Lua}
\lstset{defaultdialect=[x86masm]Assembler}
\lstset{basicstyle=\small\ttfamily,breaklines=true}
\usepackage{longtable,booktabs,array}
\usepackage{multirow}
\usepackage{calc} % for calculating minipage widths
% Correct order of tables after \paragraph or \subparagraph
\usepackage{etoolbox}
\makeatletter
\patchcmd\longtable{\par}{\if@noskipsec\mbox{}\fi\par}{}{}
\makeatother
% Allow footnotes in longtable head/foot
\IfFileExists{footnotehyper.sty}{\usepackage{footnotehyper}}{\usepackage{footnote}}
\makesavenoteenv{longtable}
\ifLuaTeX
  \usepackage{luacolor}
  \usepackage[soul]{lua-ul}
\else
  \usepackage{soul}
\fi
\setlength{\emergencystretch}{3em} % prevent overfull lines
\providecommand{\tightlist}{%
  \setlength{\itemsep}{0pt}\setlength{\parskip}{0pt}}
\setcounter{secnumdepth}{3}
\ifLuaTeX
\usepackage[bidi=basic]{babel}
\else
\usepackage[bidi=default]{babel}
\fi
\babelprovide[main,import]{english}
% get rid of language-specific shorthands (see #6817):
\let\LanguageShortHands\languageshorthands
\def\languageshorthands#1{}
\usepackage{bookmark}
\IfFileExists{xurl.sty}{\usepackage{xurl}}{} % add URL line breaks if available
\urlstyle{same}
\hypersetup{
  pdftitle={LaTeX.css},
  pdflang={en},
  pdfkeywords={latex.css,css library,class-less css,latex css},
  pdfcreator={LaTeX via pandoc}}
\usepackage{amsthm}
\newtheorem{lem}{Lemma}
\newtheorem{thm}{Theorem}
\newtheorem{dfn}{Definition}
\title{{L{a}T{e}X}.css}
\author{Vincent Dörig\\
May 2020}
\date{}



\begin{document}
\maketitle
\begin{abstract}
This almost class-less CSS library turns your HTML document into a
website that looks like a {L{a}T{e}X} document. Write semantic HTML, add
\lstinline!<link rel="stylesheet" href="https://latex.now.sh/style.css">!
to the \lstinline!<head>! of your project and you are good
to go. The source code can be found on GitHub at
\url{https://github.com/vincentdoerig/latex-css}.
\end{abstract}

{
\setcounter{tocdepth}{3}
\tableofcontents
}
\section{Getting Started}\label{getting-started}

\begin{itemize}
\item
  Add
  \lstinline!<link rel="stylesheet" href="https://latex.now.sh/style.css">!
  to the \lstinline!<head>! of your website or install the
  package using \lstinline!npm install latex.css!.
\item
  (optional) Add any classes to elements described in the
  \hyperref[class-based-elements]{next section}.
\item
  (optional) Add dark mode by adding the
  \lstinline!latex-dark-auto! class to the
  \lstinline!<body>!. More infos
  \hyperref[dark-mode]{here}.
\item
  (optional) If you need support for {L{a}T{e}X} math, add the following
  script to include \href{https://www.mathjax.org/}{MathJax}:

\begin{lstlisting}[language=HTML]
<script id="MathJax-script" async src="https://cdn.jsdelivr.net/npm/mathjax@3/es5/tex-mml-chtml.js"></script>
\end{lstlisting}
\item
  (optional) If you need syntax highlighting, add the following script
  to include \href{https://prismjs.com/}{Prism} and the Prism LaTeX
  stylesheet (or use any other):

\begin{lstlisting}[language=HTML]
<link rel="stylesheet" href="https://latex.now.sh/prism/prism.css">

<script src="https://cdn.jsdelivr.net/npm/prismjs/prism.min.js"></script>
\end{lstlisting}
\item
  Done.
\end{itemize}

\section{Class-based Elements}\label{class-based-elements}

\subsection{Author and Abstract}\label{author-abstract}

Use the following code to add an author and abstract to your document.
It will look like \hyperref[top]{this}.

\begin{lstlisting}[language=HTML]
<p class="author">John Doe <br> December 7, 2020</p>

<div class="abstract">
  <h2>Abstract</h2>
  <p>...</p>
</div>
\end{lstlisting}

\subsection{Theorems, Definitions, Lemmas and Proofs}\label{tdpl}

Theorems, definitions, lemmas and proofs are supported. Just wrap your
content in a div and add the corresponding class to the element like in
the following example.

\begin{lstlisting}[language=HTML]
<div class="theorem">...</div>
<div class="definition">...</div>
<div class="lemma">...</div>
<div class="proof">...</div>
\end{lstlisting}

Below are some examples.

\subsubsection{Proofs \& Theorems}\label{proofs-theorems}

\begin{thm}

The real numbers \(\mathbb{R}\) are uncountable.

\end{thm}

\begin{proof}

If \(\mathbb{R}\) is countable, then {[}0, 1{]} is countable as well.
Hence there exists a map C from \(\mathbb{N}\) onto {[}0, 1{]} with
\[C(n)=\sum_{i=1}^{\infty} c_{i}(n) 10^{-i}\] where \(c_{i}(n)
        \in\{0,1, \ldots, 9\},\) are the digits in decimal expansion.
Now consider a real number \[x=\sum_{i=1}^{\infty} \bar{c}_{i} 10^{-i}
        \in[0,1]\] with \(\bar{c}_{i} \neq c_{i}(i)\). Obviously
\(C(n) \neq x\) for all \(n \in \mathbb{N} .\) Hence \(C\) is not onto.
A contradiction.

\end{proof}

\begin{thm}

If \(S\) is both countable and infinite, then there is a bijection
between \(S\) and \(\boldsymbol{N}\)
itself.\textsuperscript{\phantomsection\label{ref1}\hyperref[fn1]{1}}

\end{thm}

\begin{proof}

For any \(s \in S,\) we let \(f(s)\) denote the value of \(k\) such that
\(s\) is the \(k\) the smallest element of \(S .\) This map is well
defined for any \(s,\) because there are only finitely many natural
numbers between 1 and \(s .\) It is impossible for two different
elements of \(S\) to both be the \(k\) the smallest element of \(S\).
Hence \(f\) is one-to-one. Also, since \(S\) is infinite, \(f\) is onto.

\end{proof}

\subsubsection{Lemmas}\label{lemmas}

\begin{lem}

An even number plus an even number results in an even number.

\end{lem}

\subsubsection{Definitions}\label{definitions}

\begin{dfn}

A definition is a a statement of the meaning of a word or word group or
a sign or
symbol.\textsuperscript{\phantomsection\label{ref2}\hyperref[fn2]{2}}

\end{dfn}

\subsection{Paragraphs}\label{paragraphs}

In order to get automatic first line indentation of paragraphs, like in
default {L{a}T{e}X} articles, the class
\lstinline!indent-pars! can be used with
\lstinline!article! tag.

\begin{lstlisting}[language=HTML]
<article class="indent-pars">
  ...
</article>
\end{lstlisting}

The CSS style provided follows \lstinline!babel! specific
language rules: by default, the first paragraph after a section title is
not indented in English, unlike Spanish and French languages.

To avoid first line indentation of some specific paragraph, the class
\lstinline!no-indent! can be used.

\begin{lstlisting}[language=HTML]
<p class="no-indent">...</p>
\end{lstlisting}

\subsection{Table Classes}\label{table-classes}

\subsubsection{Custom Table Borders}\label{table-borders}

To add custom borders to the table, you should use the class
\lstinline!borders-custom! with
\lstinline!<table>! tag, to begin with a table without any
border. The classes \lstinline!border-<position>-thin! and
\lstinline!border-<position>-thick!, are available, where
\lstinline!<position>! can take one of the values:
\lstinline!top!, \lstinline!right!,
\lstinline!bottom!, \lstinline!left!.

\begin{lstlisting}[language=HTML]
<table class="borders-custom ...">
  <tr class="border-bottom-thick">
    <td>...</td><td>...</td><td>...</td>
  </tr>
  <tr class="border-bottom-thin">
    <td>...</td><td>...</td><td>...</td>
  </tr>
  ...
\end{lstlisting}

See full examples in the \hyperref[tables]{section about tables}.

\subsubsection{Table Column Alignment}\label{table-column-aligment}

For each column of the table, there is one class for left, center or
right alignment, up to 12 columns: \lstinline!col-<n>-l!,
\lstinline!col-<n>-c!,
\lstinline!col-<n>-r!, \lstinline!n!=1,
...,12. For example, a table with 3 columns using the following classes

\begin{lstlisting}[language=HTML]
<table class="col-1-l col-2-c col-3-r">
  <tr>
    <td>...</td><td>...</td><td>...</td>
  </tr>
  ...
\end{lstlisting}

aligns the first column to the left, centers the second one and aligns
the third to the right.

See full examples in the \hyperref[tables]{section about tables}.

\section{Language Support}\label{language-support}

The labels of theorems, definitions, lemmas and proofs can be changed to
other
\href{https://github.com/vincentdoerig/latex-css/tree/master/lang}{supported
language} by including the following snippet, linking to the desired
language in addition to the main CSS file.

\begin{lstlisting}[language=HTML]
<link rel="stylesheet" href="https://latex.now.sh/lang/es.css" />
  ...
<html lang="es">
\end{lstlisting}

Take a peek at the \href{/languages}{language support demo} to see how
the labels of the different languages change.

\section{Sidenotes}\label{sidenotes}

Sidenotes can be used as an alternative to footnotes, where the user
does not have to jump to the bottom of the page to read it. On mobile,
click the superscript to reveal the notek\marginpar{\small {Yay,
sidenotes!. If you are on mobile, I will appear inline. If you are using
a larger screen, the sidenote will appear to the right of the text.}}.

Sidenotes do need a little bit of setup, they are made up of a label, an
invisible checkbox on top of the number and a span with the text inside.
The superscript is set automatically and incremented using CSS when the
checkbox has a class of \lstinline!sidenote-number!.

\begin{lstlisting}[language=HTML]
<label for="sn-1" class="sidenote-toggle sidenote-number"></label>
<input type="checkbox" id="sn-1" class="sidenote-toggle" />
<span class="sidenote"><!-- sidenote content --></span>
\end{lstlisting}

If you do not need superscripted numbers, you can opt out of the
\lstinline!sidenote-number! class and the sidenote will
not have a number assigned. On a smaller screen, you will need to add a
symbol inside the \lstinline!label! for the user to click
on. \marginpar{\small {This is a sidenote without a number.}}

The snippet for a sidenote without a number is very similar:

\begin{lstlisting}[language=HTML]
<label for="sn-anything" class="sidenote-toggle">*</label>
<input type="checkbox" id="sn-anything" class="sidenote-toggle" />
<span class="sidenote"><!-- sidenote content --></span>
\end{lstlisting}

Add a class of \lstinline!left! to the span with the
\lstinline!sidenote! class to make the note appear on the
left side of the page on instead of right.

The symbol you could use to indicate a sidenote is up to you. What about
this notebook. \marginpar{\small {A notebook indicating a note. Aha.\\
(if on a large screen, resize to mobile to see the emoji)}}

\section{HTML Elements}\label{html-elements}

For a preview of all HTML elements with LaTeX.css, check out the
\href{/elements}{HTML5 elements test page}.

\subsection{Text Formatting}\label{text-formatting}

This sentence is \textbf{bold}. If you like semantics, you might go with
\textbf{strong} or \emph{emphasized} text. If not, \emph{italic} is
still around. {Small} text is for fine print. Your copy can also be
\textsubscript{subscripted} and \textsuperscript{superscripted},
\ul{inserted}, \st{deleted}, or \hl{highlighted}. You would use a
\hyperref[ux21]{hyperlink} to go to a new page. Keyboard input elements
like {Cmd} + {Shift} are used to display textual user input.

\subsection{Blockquotes}\label{blockquotes}

\begin{quote}
Give me six hours to chop down a tree and I will spend the first four
sharpening the axe. --- Abraham Lincoln
\end{quote}

\subsection{Definition Lists}\label{definition-lists}

\begin{description}
\tightlist
\item[Definition Title One] First definition description
\item[Binomial theorem]

\[(x+y)^{n}=\sum_{k=0}^{n}\left(\begin{array}{l}n \\
k\end{array}\right) x^{n-k}
y^{k}=\sum_{k=0}^{n}\left(\begin{array}{l}n \\ k\end{array}\right)
x^{k} y^{n-k}\]
          
\end{description}

\subsection{Tables}\label{tables}

% Tables have been simplified, because Pandoc LaTeX reader and writer don't
% support all features from `longtable` and `multirow` packages.
% Refs:
% - https://github.com/jgm/pandoc/issues/6311
% - https://github.com/jgm/pandoc/issues/6596

\begin{longtable}[]{@{}lll@{}}
\caption{A sample table with a descriptive caption.}\tabularnewline
\toprule\noalign{}
Header 1 & Header 2 & Header 3 \\
\midrule\noalign{}
\endhead
Description 1 & Description 2 & Description 3 \\
Description 1 & Description 2 & Description 3 \\
Description 1 & Description 2 & Description 3 \\
\midrule\noalign{}
Footer 1 & Footer 2 & Footer 3 \\
\bottomrule\noalign{}
\end{longtable}

\begin{longtable}[]{@{}lll@{}}
\caption{Example table taken from
\href{https://texdoc.org/serve/booktabs/0}{this paper} on how to produce
quality tables with LaTeX.}\tabularnewline
\toprule\noalign{}
Animal & Description & Price (\$) \\
\midrule\noalign{}
\endhead
\bottomrule\noalign{}
\endlastfoot
Gnat & per gram & 13.65 \\
Gnu & stuffed & 92.50 \\
Emu & stuffed & 33.33 \\
Armadillo & frozen & 8.99 \\
\end{longtable}

\begin{longtable}[]{@{}llll@{}}
\caption{Color names and values.}\tabularnewline
\toprule\noalign{}
Name & HEX & HSL & RGB \\
\midrule\noalign{}
\endhead
\bottomrule\noalign{}
\endlastfoot
Teal & \texttt{\#008080} & \texttt{hsl(180,\ 100\%,\ 25\%)} &
\texttt{rgb(0,\ 128,\ 128)} \\
Goldenrod & \texttt{\#daa520} & \texttt{hsl(43,\ 74\%,\ 49\%)} &
\texttt{rgb(218,\ 165,\ 32)} \\
Cornflowerblue & \texttt{\#6495ed} & \texttt{hsl(219,\ 79\%,\ 66\%)} &
\texttt{rgb(100,\ 149,\ 237)} \\
Lightcoral & \texttt{\#f08080} & \texttt{hsl(0,\ 79\%,\ 72\%)} &
\texttt{rgb(240,\ 128,\ 128)} \\
\end{longtable}

\begin{longtable}[]{@{}lllllll@{}}
\caption{Example table with custom borders and aligment (De Morgan\textquotesingle s Law:
\(\lnot P \lor \lnot Q \iff \lnot (P \land Q)\)).}\tabularnewline
\toprule\noalign{}
\(P\) & \(Q\) & \(\lnot P\) & \(\lnot Q\) & \(P \land Q\) &
\(\lnot P \lor \lnot Q\) & \(\lnot (P \land Q)\)\\
\midrule\noalign{}
\\% an empty row is needed because multirow issues
\multirow{2}{*}{T} & T & \multirow{2}{*}{F} & F & T & \textbf{F} &
\textbf{F} \\
& F & & T & F & \textbf{T} & \textbf{T} \\
\multirow{2}{*}{F} & T & \multirow{2}{*}{T} & F & F & \textbf{T} &
\textbf{T} \\
& F & & T & F & \textbf{T} & \textbf{T} \\
\bottomrule\noalign{}
\end{longtable}

\begin{longtable}[]{@{}llll@{}}
\caption{Close approaches to the Earth by NEOs.}\tabularnewline
\toprule\noalign{}
CA Date & Object Name & Diameter & Dist. Min. (LD) \\
\midrule\noalign{}
2023-Jul-10 & 2018 NW & 7.3m-16m & 0.287 \\
2023-Jul-10 & 2023 LN1 & 45m-00m & 17.813 \\
2023-Jul-11 & 2023 MD2 & 36m-80m & 5.570 \\
2023-Jul-11 & 2023 NE & 32m-70m & 11.253 \\
2023-Jul-11 & 2023 MQ1 & 28m-62m & 10.716 \\
2023-Jul-12 & 2023 OM & 20m-46m & 7.628 \\
2023-Jul-12 & 2018 UY & 190m-420m & 7.412 \\
\bottomrule\noalign{}
\end{longtable}

\subsection{Images}\label{images}

\begin{figure}[!htp]
\centering
{[}{[}Insert here the remote image:https://images.unsplash.com/photo-1477346611705-65d1883cee1e?auto=format\&fit=crop\&w=1000\&q=80{]}{]}
\caption{Mountain landscape by \href{https://unsplash.com/@heytowner}{John Towner}.}
\end{figure}

\section{Miscellaneous}\label{miscellaneous}

\subsection{Syntax Highlighting}\label{syntax-highlighting}

If you need syntax highlighting for code, LaTeX.css provides a
\href{https://prismjs.com/}{PrismJS} theme that immitates the
\href{https://github.com/gpoore/minted}{minted} package for LaTeX. Add
the following stylesheet and script:

\begin{lstlisting}[language=HTML]
<link rel="stylesheet" href="https://latex.now.sh/prism/prism.css">

<script src="https://cdn.jsdelivr.net/npm/prismjs/prism.min.js"></script>
\end{lstlisting}

And use it like this:

\begin{lstlisting}[language=HTML]
<pre>
  <code class="language-html">
    <!-- your HTML code snippet -->
  </code>
</pre>
\end{lstlisting}

Change which languages Prism highlights by customising the script
\href{https://prismjs.com/download.html}{here}.

\vspace{5pt}\rule{0.5\linewidth}{0.5pt}

1. From \url{https://www.math.brown.edu/~res/MFS/handout8.pdf}.
\hyperref[ref1]{\^{}}

2. ``Definition.'' Merriam-Webster.com Dictionary, Merriam-Webster,
\url{https://www.merriam-webster.com/dictionary/definition}. Accessed 18
May. 2020. \hyperref[ref2]{\^{}}

\end{document}
